\chapter{Problem Analysis}
\label{sec:correlated_concepts}
This chapter identifies the limitation of the \emph{Causal Concept Faithfulness Metric} addressed by this thesis. It contrasts the metric's causal assumptions with the practical construction of counterfactual questions and then shows how dependencies between concepts give rise to distinct interaction effects that single-concept interventions cannot reliably capture.

\section{Assumptions in the Observational Graph}
\label{sec:observational_graph}

In the appendix of their work, Matton et al.\ \cite{mattonWalkTalkMeasuring2024} discuss the data generating process underlying a dataset question via a causal graph (Figure~\ref{fig:matton_dgp}). The variables in this graph are interpreted as follows: The possible questions $X$ are determined by an unobserved \enquote{state of the world} variable $U$ through the mediator set $\{C_1,\dots,C_M, V\}$. Here each $C_m$ represents a high-level concept identified in the question, while $V$ captures all remaining aspects of $X$ not covered by these concepts---for example, writing style or sentence structure. $Y$ represents the answer given by model $\mathcal{M}$ and is modeled as a categorical random variable influenced by $C_1,\dots,C_M$ and $V$, while $\mathcal{E}$ denotes unobserved stochasticity of the LLM.

\begin{figure}[ht]
	\centering
	\includestandalone[mode=buildnew, width=10cm]{graphs/matton_dgp}
	\caption{The data generating process for original dataset questions and model responses, as assumed by Matton et al.\ \cite{mattonWalkTalkMeasuring2024}.}
	\label{fig:matton_dgp}
\end{figure}

Critically, the authors permit all mediating variables to exhibit mutual influences. Concepts $C_1,\dots,C_M$ may affect one another and may also influence the rest of the question text $V$. Moreover, the confounder $U$ can introduce correlations among the mediators. To indicate this layer of mutual dependence, the mediators are enclosed by a dotted boundary in Figure~\ref{fig:matton_dgp}.

On a practical level, this means that in the observational data (the given questions) concepts are often correlated and not independent. For instance, a person's name may causally influence their gender or the confounder $U$ may cause age and experience level to appear jointly within a question.

\section{Disentanglement Assumption for Interventions}

Matton et al.\ \cite{mattonWalkTalkMeasuring2024} argue that measuring the \emph{total causal effect} of a concept $C_m$ on the model's answer $Y$ would be undesirable. An explanation generated by the LLM is not expected to mention \enquote{concepts that influenced other concepts that then influenced its answer,} but rather to directly name the concepts that influenced its own prediction \cite{mattonWalkTalkMeasuring2024}. Therefore, to align the Causal Concept Effect with the Explanation-Implied Effect, the \emph{direct effect} of $C_m$ on the answer variable $Y$ must be isolated—i.e., any influence mediated by other variables should be discarded \cite{pearlCausality2009}.

Measuring the direct effect requires a \emph{surgical intervention} on a single concept $C_m$ while keeping all other concepts $C_i$ ($i \neq m$) and the rest of the question text $V$ fixed at their original values \cite{pearlCausality2009}. This operation severs any incoming dependencies to the intervened variable, rendering it independent of other causes \cite{pearlCausality2009, eberhardtInterventionsCausalInference2007}. The assumed observational correlations between the mediators do not contradict the possibility of such an intervention: structural causal models allow interventions on individual variables even when those variables are statistically associated in observational data \cite{pearlCausality2009}.

The effect of a surgical intervention on the causal graph is illustrated in Figure~\ref{fig:matton_dgp_intervened}. The value of $C_m$ is replaced by an alternative $c_m'$, while all other mediators are held constant. Because the fixed mediators are no longer influenced by the confounder $U$, the variable $U$ is removed from the graph. The disappearance of all mutual influences among the mediators is indicated by the removal of the dotted boundary.

\begin{figure}[ht]
	\centering
	\includestandalone[mode=buildnew, width=16cm]{graphs/matton_dgp_intervened}
	\caption{Causal graph under a surgical intervention on concept $C_m$; all other concepts and $V$ remain fixed.}
	\label{fig:matton_dgp_intervened}
\end{figure}

In summary, the \emph{disentanglement} assumption—that each concept can be modified independently without altering the other concepts or the rest of the question text—is introduced to enable the estimation of direct effects. It does not require the concepts to be independent in the original data; it only demands that the intervention itself be feasible in a way that leaves everything else unchanged.

\section{The Practical Limit}
\label{sec:pratical_limit}
While disentanglement is theoretically coherent, several practical obstacles arise. First, it may be idealistic to assume that $V$ can be kept perfectly fixed under every intervention; certain concept changes might result in unnatural or incoherent phrasing if the surrounding text is left untouched.

More importantly, in many real-world datasets concepts are correlated in the training data of the LLM. Matton et al.\ \cite{mattonWalkTalkMeasuring2024} themselves acknowledge that when two concepts are strongly associated, intervening on one while holding the other constant can produce unreliable direct-effect estimates. This problem becomes especially acute for removal-based counterfactuals: the LLM may be able to reconstruct the removed concept value from the remaining, unchanged concepts because it has learned to expect them to co-occur. For example, in a job application setting, an applicant’s gender may be easily inferred from their first name even if gender is explicitly removed.

Matton et al.\ \cite{mattonWalkTalkMeasuring2024} illustrate the issue with a question from the MedQA dataset. They observe that the concepts \emph{the patient’s eating disorder} and \emph{the patient’s self-perception of weight} each receive a surprisingly low Causal Concept Effect. They hypothesize that this occurs because the LLM reconstructs one concept from the other when only one is removed, effectively turning the intended surgical intervention into a non-surgical modification that leaves traces in the input.

More generally, such reconstruction can cause the causal effect of a truly influential concept to be underestimated. This, in turn, may artificially widen the gap between the Causal Concept Effect and the Explanation-Implied Effect, potentially penalizing an otherwise faithful LLM with an unfairly low faithfulness score. Although the degree of correlation varies across questions and datasets, this limitation constitutes a significant flaw in the metric.
%%continue review here
\section{Interaction between Concepts}
\label{sec:interaction_effects}
The example by Matton et al.\ \cite{mattonWalkTalkMeasuring2024}, in which the model reconstructs a removed concept from another remaining concept, represents a redundancy-like interaction with respect to the model's answer. More generally, concepts can affect a response additively, redundantly, or synergistically, expanding the \textit{Correlated Concept Problem} beyond the feasibility of surgical interventions. This terminology is motivated by information theory and, in particular, \emph{Partial Information Decomposition} (PID), which decomposes how multiple source variables jointly provide information about a target variable \cite{williamsNonnegativeDecompositionMultivariate2010, kolchinskyNovelApproachPartial2022}. PID is normally defined using information-theoretic quantities, whereas the function $J_{\mathbf{x}}$ introduced below is based on changes in the evaluated model's answer distribution. The three information-theoretic categories are therefore mapped here to additive, subadditive, and superadditive relationships between intervention effects; they are not claimed to constitute a formal PID of the model output. Related information-sharing categories have also been incorporated into causal frameworks \cite{martinez-sanchezDecomposingCausalityIts2024}.

For illustration, consider an initially masked question $\mathbf{x}$ into which two concepts $C_i$ and $C_j$ are introduced. Throughout this section only, the values $c_i$ and $c_j$ denote the masked (i.e., absent) state of the corresponding concepts, whereas $c_i'$ and $c_j'$ denote arbitrary non-null values. For the conceptual classification developed here, let $J_{\mathbf{x}}(T)$ denote the effect on the evaluated model's answer distribution of simultaneously intervening on every concept in a non-empty set $T$. For a singleton set, this notation coincides with the Causal Concept Effect introduced in Section~\ref{sec:problem_setting}:

\begin{equation}
	J_{\mathbf{x}}(\{C_m\}) := \mathrm{CE}(\mathbf{x},C_m).
\end{equation}

Thus, $J_{\mathbf{x}}(\{C_i,C_j\})$ denotes the joint effect obtained by replacing $(c_i,c_j)$ with $(c_i',c_j')$ simultaneously, while $J_{\mathbf{x}}(\{C_i\})$ and $J_{\mathbf{x}}(\{C_j\})$ denote the corresponding individual effects. At this point, $J_{\mathbf{x}}$ is an abstract joint-intervention effect used only to describe response-level interaction regimes. The formal construction of joint interventions is deferred to Section~\ref{sec:joint_effect}, and FELICS's corresponding question-specific value function is defined later in Equation~\ref{eq:value_function}. Depending on the relationship between these quantities, three regimes can be distinguished.

\subsection{Additivity}
Under \emph{additivity}, the joint intervention effect equals the sum of the individual intervention effects,
\begin{equation}
	J_{\mathbf{x}}(\{C_i,C_j\}) = J_{\mathbf{x}}(\{C_i\})+J_{\mathbf{x}}(\{C_j\}).
\end{equation}
This equality expresses an absence of interaction on the chosen effect scale. It does not imply that $C_i$ and $C_j$ are statistically independent in the data; associated inputs may have additive effects on the response, and statistically independent inputs may nevertheless exhibit a non-additive response-level interaction.

As a simplified credit-scoring example, the number of recent credit inquiries and the requested loan duration could contribute approximately independent evidence.

\subsection{Redundancy}
In the mapped, response-level sense used here, concepts exhibit \emph{redundancy-like} behavior when their individual effects overlap. Introducing information that is already provided by the other concept has only a limited additional effect, so the combined intervention is smaller than the sum of the individual interventions,
\begin{equation}
	J_{\mathbf{x}}(\{C_i,C_j\}) < J_{\mathbf{x}}(\{C_i\})+J_{\mathbf{x}}(\{C_j\}).
\end{equation}

In the credit-scoring example, homeownership and the absence of monthly rent may provide redundant information about housing expenses.

\subsection{Synergy}
In a \emph{synergy-like} interaction, an additional effect emerges only when both concepts are considered jointly. Consequently, the combined intervention produces a larger effect than expected from the sum of the individual interventions,
\begin{equation}
	J_{\mathbf{x}}(\{C_i,C_j\}) > J_{\mathbf{x}}(\{C_i\})+J_{\mathbf{x}}(\{C_j\}).
\end{equation}

High income and low existing debt may exhibit synergy within the credit-scoring example when their combination provides stronger evidence of repayment capacity than either property provides separately.

\subsection{Implications for Causal Concept Effects}
The examples above are restricted to the interaction of two concepts for simplicity. In practice, multiple concepts may interact simultaneously, giving rise to substantially more complex information-sharing structures. Williams and Beer \cite{williamsNonnegativeDecompositionMultivariate2010} generalize Partial Information Decomposition to arbitrary numbers of variables and introduce corresponding partial information diagrams to represent these relationships.

Moreover, the discussed examples describe the effect of \emph{introducing} concepts into an initially masked question. This formulation was chosen because it allows the three response-level regimes to be interpreted intuitively in terms of the additional effect contributed by a concept. Matton et al.\ \cite{mattonWalkTalkMeasuring2024}, however, estimate causal concept effects by \emph{removing} concepts from the original question. Under this complementary perspective, the inequalities for redundancy-like and synergy-like behavior are reversed. Removing two concepts with redundant information jointly eliminates more information than removing either concept individually, yielding

\begin{equation}
	J_{\mathbf{x}}(\{C_i,C_j\})
	>
	J_{\mathbf{x}}(\{C_i\})+J_{\mathbf{x}}(\{C_j\}),
\end{equation}

whereas removing two synergistic concepts results in

\begin{equation}
	J_{\mathbf{x}}(\{C_i,C_j\})
	<
	J_{\mathbf{x}}(\{C_i\})+J_{\mathbf{x}}(\{C_j\}).
\end{equation}

Consequently, the MedQA example discussed by Matton et al. \cite{mattonWalkTalkMeasuring2024}, in which one removed concept can be reconstructed from another, is consistent with redundancy-like behavior under the removal interpretation.

Statistical association and response-level interaction must be distinguished. Association concerns the joint distribution of the concepts: observing one concept may make the value of another more predictable. Interaction, as used in this thesis, concerns their effects on a response variable: the effect of intervening on one concept may depend on whether or how another concept is intervened upon, with the evaluated model's answer distribution serving as the response. Strong association between concepts can indicate the possibility of redundant information and may therefore serve as a useful screening proxy for redundancy-like effects. It does not, however, establish redundancy or interaction with respect to the model output. Conversely, response-level interaction---particularly synergy---may occur without observable pairwise association between the concepts \cite{picaInvariantComponentsSynergy2017}.

To recapitulate, the original \emph{Causal Concept Faithfulness Metric} assumes that each concept can be intervened upon surgically while the remaining input is held fixed and evaluates only singleton intervention effects. Redundancy-like relationships can undermine the intended intervention because the removed information may be reconstructed from another concept, while synergy-like effects may become observable only when multiple concepts are modified jointly. Consequently, singleton interventions may not reliably characterize causal concept effects when such response-level interactions are present.

With the correlated concepts problem established and the response-level interaction regimes operationalized, the remainder of this thesis focuses on constructing and validating an extension of the \emph{Causal Concept Faithfulness Metric} that evaluates plausible concept interactions during the estimation of causal concept effects.
